\documentclass[10pt,a4paper]{article}
\usepackage{amsmath,amssymb,amsthm,graphicx}
\usepackage{amsmath,amssymb,amsthm,algorithmicx,algpseudocode}
\usepackage[no-math]{fontspec}
\newenvironment{authorclaim}[1]{\par\medskip\noindent\textbf{#1.}\ \ignorespaces}{\par\medskip}
\allowdisplaybreaks
\title{Exponential tree-like semantic Res(xor) refutation lower bound for perfect matching}
\author{Dmitry Itsykson, Artur Riazanov}
\date{doi:10.4230/LIPIcs.CCC.2021.3; CC BY 4.0}
\begin{document}
\maketitle
\noindent\textit{Editorial scope.} Theorem2 is counted once with its full proof-to-communication bridge and complete AppendixA. Supporting Lemmas1/12/17 and direct matrix lower bound are not split into counted tasks.\par\medskip
\noindent\textit{Allowed standard prerequisites (editorial).}\begin{itemize}
\item Randomized communication complexity, unique/set disjointness lower bounds and biased random-walk hitting time. Formal Section2 Theorems7–8 and AppendixA Theorem45, with original full33-entry bibliography. Standard imported results retained as cited prerequisites, not recursively reproved.
\end{itemize}\medskip
% Faithful typesetting transcription of necessary formal Section2 definitions.
\section*{2 Preliminaries}
\paragraph{Notations.} We use the following notation: $[n]=\{1,2,\ldots,n\}$. Let $S^{n\times m}$ denote the set of matrices of size $n\times m$ with elements from $S$. We denote by $0_{n\times m}$ the zero matrix of size $n\times m$ and by $1_{n\times m}$ the matrix of the same size containing only ones. For square matrices $A_1,\ldots,A_k$ we denote a diagonal block matrix with blocks $A_1,\ldots,A_k$ by $\operatorname{diag}(A_1,\ldots,A_k)$. For $x\in\{0,1,\ldots,2^k-1\}$ we denote a vector $(a_0,\ldots,a_{k-1})\in\{0,1\}^k$ such that $x=\sum_{i=0}^{k-1}a_i2^i$ by $\operatorname{bin}_k(x)$, i.e. $(a_0,\ldots,a_{k-1})$ is the reversed binary representation of $x$. For vectors $v_1,\ldots,v_n$ from a vector space over a field $\mathbb F$ we denote their linear span by $\operatorname{Span}(v_1,\ldots,v_n)$. We use coordinate-wise comparison of strings from $\{0,1\}^n$, i.e. for $x,y\in\{0,1\}^n$ we write $x\le y$ iff $x_i\le y_i$ for each $i\in[n]$. We denote the set of variables of a CNF-formula $\phi$ by $\operatorname{Vars}(\phi)$.
\paragraph{Communication complexity.} We briefly recall some notions of communication complexity. For formal definition and details we refer to [24].
In the classic two-party randomized communication protocol with public randomness, Alice and Bob cooperate to compute a relation $Q\subseteq X\times Y\times Z$: Alice has an input $x\in X$ and Bob has an input $y\in Y$, their goal is to compute $z\in Z$ such that $(x,y,z)\in Q$. We assume that Alice and Bob have access to an arbitrary large random string of bits that is common for Alice and Bob. Let for every $x\in X$ and $y\in Y$, $R^\delta_{\mathrm{pub}}(Q,x,y)$ denote the minimal number of bits Alice and Bob need to transmit between each other so they both find a $z\in Z$ such that $(x,y,z)\in Q$ with probability at least $1-\delta$ taken over the values of the common random string. And $R^\delta_{\mathrm{pub}}(Q):=\max_{x\in X,y\in Y}R^\delta_{\mathrm{pub}}(Q,x,y)$.
We also consider multiparty communication protocols in the number on forehead (NOF) model that extends two-party protocols for an arbitrary number of parties. In this setting $k$ parties cooperate to compute a relation $Q\subseteq X_1\times X_2\times\cdots\times X_k\times Y$. The $i$th party has $x_i\in X_i$ written on their forehead so they know all $x_j$ for $j\ne i$, their goal is to compute $y\in Y$ such that $(x_1,x_2,\ldots,x_k,y)\in Q$. The parties communicate by taking turns broadcasting messages to all other parties until all parties learn the value of $y\in Y$ such that $(x_1,\ldots,x_k,y)\in Q$. In this model we also assume that all parties have access to a common random string of bits. Let $R^\delta_{\mathrm{pub}}(Q,x_1,\ldots,x_k)$ for $x_1\in X_1,\ldots,x_k\in X_k$ denote the minimal total number of bits transmitted until each party learns $y\in Y$ such that $(x_1,\ldots,x_k,y)\in Q$ with probability at least $1-\delta$ taken over the set of values of the random string of bits. Also, let $R^\delta_{\mathrm{pub}}(Q):=\max_{x_1\in X_1,\ldots,x_k\in X_k}R^\delta_{\mathrm{pub}}(Q,x_1,\ldots,x_k)$.
Let $f$ be a function from $X_1\times X_2\times\cdots\times X_k\to Y$. Then $R^\delta_{\mathrm{pub}}(f)$ denotes $R^\delta_{\mathrm{pub}}(Q_f)$, where $Q_f=\{(x_1,x_2,\ldots,x_k,y)\mid f(x_1,\ldots,x_k)=y\}$.
We prove communication complexity lower bounds by reduction from different versions of the set disjointness problem. $\mathrm{DISJ}_{k,n}$ is a function $\{0,1\}^{kn}\to\{0,1\}$ such that for every $x_1,\ldots,x_k\in\{0,1\}^n$ the following holds:
\[\mathrm{DISJ}_{k,n}(x_1,\ldots,x_k)=\bigwedge_{j=1}^n\neg\left(\bigwedge_{i=1}^k(x_i)_j\right).\]
Let us define the communication promise problem $\mathrm{UDISJ}_{k,n}$ in the $k$-party NOF model. For each $i\in[k]$ the string $x_i$ is written on the forehead of the $i$th party, it is guaranteed that there exists at most one index $j\in[n]$ such that for every $i\in[k]$, $(x_i)_j=1$. The goal is to compute $\mathrm{DISJ}_{k,n}(x_1,\ldots,x_k)$.
\begin{authorclaim}{Theorem7 ([31,32])}
$R^{1/3}_{\mathrm{pub}}(\mathrm{UDISJ}_{k,n})=\Omega(\sqrt n/(2^kk))$.
\end{authorclaim}
For $k=2$ we omit the first index: $\mathrm{DISJ}_n=\mathrm{DISJ}_{2,n}$; in this case Theorem7 may be improved.
\begin{authorclaim}{Theorem8 ([16])}
$R^{1/3}_{\mathrm{pub}}(\mathrm{DISJ}_n)\ge R^{1/3}_{\mathrm{pub}}(\mathrm{UDISJ}_{2,n})=\Omega(n)$.
\end{authorclaim}
\paragraph{Proof complexity.} We consider refutational proof systems for the language of unsatisfiable CNF-formulas UNSAT. A refutation of $\phi\in\mathrm{UNSAT}$ in a proof system $\Pi$ is a sequence of Boolean functions (proof lines) such that each proof line either represents a clause of $\phi$ or derived from previous proof lines in the sequence via some sound inference rules. The last line of the proof is identically zero function. A proof system $\Pi$ is defined by a representation of proof lines and by a set of admissible inference rules. It is required that the inference rules are polynomially verifiable i.e. there exists an algorithm that checks whether it is legitimate to derive a line $L_0$ from the lines $L_1,\ldots,L_k$.
For example, in the Resolution proof lines are represented by clauses and the only inference rule is the resolution rule that allows deriving a clause $A\vee B$ from the clauses $A\vee x$ and $A\vee\neg x$.
The size of a proof is the total size of all representations of proof lines in the proof. The length of a proof is the number of proof lines in it.
A tree-like proof is such a proof that every its line can be used as a premise of a rule at most once. For each proof system, we can also consider its tree-like version where all proofs are constrained to be tree-like.
We also consider semantic refutational proof systems, where we drop the requirement for polynomial verification of inference rules i.e. we allow to derive any sound consequence from the premises. For such systems it is crucial to bound fan-in i.e. the number of the premises from which each proof line can be derived, otherwise, it would be possible to derive a contradiction from the clauses of the initial formula immediately. For example, it is well-known that Resolution is polynomially equivalent to a semantic proof system with fan-in 2 operating with clauses.
A lower bound on the proof size in a semantic proof system implies a lower bound on the proof size in its syntactic counterpart because a syntactic proof is always a semantic proof that operates with the same class of proof lines.
We define semantic $\mathrm{Res}(\oplus)$ as a semantic proof system with fan-in 2 that operates with linear clauses. A linear clause is a disjunction of linear equations over $\mathbb F_2$: $\bigvee_{i=1}^k(f_i=a_i)$, where $f_i$ is a linear form over $\mathbb F_2$ and $a_i\in\mathbb F_2$. Notice that an ordinary clause $\bigvee_{i\in P}x_i\vee\bigvee_{j\in N}\neg x_j$ can be represented by the linear clause $\bigvee_{i\in P}(x_i=1)\vee\bigvee_{j\in N}(x_j=0)$. For definition of syntactic version of $\mathrm{Res}(\oplus)$ we refer to [15]; it is also proved there that syntactic and semantic $\mathrm{Res}(\oplus)$ are polynomially equivalent.
We define semantic $\mathrm{Res}(\mathrm{PC}_d)$ as a semantic proof system with fan-in 2 that operates with disjunctions of equations of type $f=0$, where $f$ is a degree-$d$ polynomial over $\mathbb F_2$. Notice that semantic $\mathrm{Res}(\mathrm{PC}_1)$ is exactly semantic $\mathrm{Res}(\oplus)$. For the definition of the syntactic version of $\mathrm{Res}(\mathrm{PC}_d)$ we refer to [19].
Following [1] we define $\mathrm{Th}(k)$ as a semantic proof system with fan-in 2 that operates with polynomial inequalities $g\ge0$, where $g$ is a polynomial of degree at most $k$ with integer coefficients and Boolean variables. A clause $\bigvee_{i\in P}x_i\vee\bigvee_{j\in N}\neg x_j$ can be represented by an inequality $\sum_{i\in P}x_i+\sum_{j\in N}(1-x_j)-1\ge0$.
\paragraph{Proof complexity and communication complexity.} For an unsatisfiable CNF-formula $\phi$ we define the communication problem $\operatorname{Search}(\phi)$. Search$(\phi)$ is the following problem: given an assignment of the variables of the unsatisfiable CNF $\phi$, find a clause that is falsified by this assignment. It is assumed that variables of $\phi$ are somehow partitioned between the parties.
Following the paper [9] we consider a semantic proof system $\mathrm T^{\mathrm{cc}}(k,c)$ that models many interesting syntactic and semantic proof systems. The proof lines in $\mathrm T^{\mathrm{cc}}(k,c)$ can be arbitrary Boolean functions having the following property: for every proof line $C$ and every partition of variables of $C$ between $k$ parties, the NOF $k$-party randomized communication complexity of $C$ is at most $c$ w.r.t. this partition. We also define a semantic proof system $\mathrm T^{\mathrm{cc}}_{\mathrm{os}}(k,c)$ that is a subsystem of $\mathrm T^{\mathrm{cc}}(k,c)$ with the restriction that a communication protocol for proof lines must have a one-sided error: if the value of a proof line is zero, then the protocol should return zero with probability 1.
For example, $\mathrm T^{\mathrm{cc}}(2,2)$ simulates Resolution; $\mathrm T^{\mathrm{cc}}(2,O(1))$ simulates $\mathrm{Res}(\oplus)$ [22,15]; $\mathrm T^{\mathrm{cc}}(k,O(k^3\log^2n))$, where $n$ is the number of variables in a refuted formula, simulates $\mathrm{Th}(k-1)$ [9]. In Section3 we show that $\mathrm T^{\mathrm{cc}}_{\mathrm{os}}(d+1,O(1))$ simulates $\mathrm{Res}(\mathrm{PC}_d)$.
The following connection between the communication complexity of Search$(\phi)$ and tree-like proof complexity of $\phi$ is known.
\begin{authorclaim}{Lemma9 ([1,9])}
If a CNF formula $\phi$ has a tree-like $\mathrm T^{\mathrm{cc}}(k,c)$ refutation of length $\ell$ then, over any $k$-partition of the variables, there is a randomized bounded-error $k$-party NOF protocol for Search$(\phi)$ with communication cost $O(c\cdot\log\ell\log\log\ell)$.
\end{authorclaim}
In Section3 we show that for $\mathrm T^{\mathrm{cc}}_{\mathrm{os}}(k,c)$ the bound can be improved, see Remark14.
\paragraph{Basic formulas.} A CNF formula $\mathrm{PHP}_n^m$ encodes the pigeonhole principle; $\mathrm{PHP}_n^m$ states that it is possible to put $m$ pigeons into $n$ holes such that every pigeon flies to at least one hole and at most one pigeon flies to each hole. $\mathrm{PHP}_n^m$ depends on variables $p_{i,j}$ for $i\in[m]$ and $j\in[n]$ and $p_{i,j}=1$ iff the $i$-th pigeon flies to the $j$-th hole. $\mathrm{PHP}_n^m$ is the conjunction of $m(m-1)n/2$ hole axioms and $m$ pigeons axioms. For every $i\in[m]$ $\mathrm{PHP}_n^m$ contains a pigeon axiom $(p_{i,1}\vee p_{i,1}\vee\cdots\vee p_{i,n})$. And for every $j\in[n]$ and every $k\ne\ell\in[n]$, $\mathrm{PHP}_n^m$ contains a hole axiom $(\neg p_{k,j}\vee\neg p_{\ell,j})$. $\mathrm{PHP}_n^m$ is unsatisfiable iff $m>n$.
For an undirected graph $G(V,E)$, the formula $\mathrm{PM}_G$ encodes in CNF that $G$ has a perfect matching. The formula $\mathrm{PM}_G$ has $|E|$ variables, each of them corresponds to an edge of $G$, $x_e$ is the variable corresponding to $e\in E$.
\[\mathrm{PM}_G=\bigwedge_{v\in V}\left(\left(\bigvee_{e\text{ is incident to }v}x_e\right)\wedge\bigwedge_{e_1\ne e_2\text{ are incident to }v}(\neg x_{e_1}\vee\neg x_{e_2})\right).\]
$\mathrm{PM}_G$ is unsatisfiable iff $G$ does not have a perfect matching.
\begin{authorclaim}{Theorem10 ([26])}
Let $G$ be a graph with $n$ vertices, which has no perfect matching. Then the formula $\mathrm{PM}_G$ has a tree-like $\mathrm{Res}(\oplus)$ refutation of size $2^{O(n)}$.
\end{authorclaim}
\begin{authorclaim}{Proposition11 ([14])}
Let $G$ be a graph with an odd number of vertices. Then the formula $\mathrm{PM}_G$ has a tree-like $\mathrm{Res}(\oplus)$ refutation of size $\operatorname{poly}(n)$.
\end{authorclaim}
The binary pigeonhole principle $\mathrm{BPHP}_{2^\ell}^m$ states that there are $m$ different $\ell$-bit binary strings $s_1,s_2,\ldots,s_m$. $\mathrm{BPHP}_{2^\ell}^m$ has $m\ell$ variables corresponding to the bits of $s_i$ for $i\in[m]$. Then $\mathrm{BPHP}_{2^\ell}^m=\bigwedge_{i\ne j\in[m]}s_i\ne s_j$, where the predicate $s_i\ne s_j$ is encoded as a $2\ell$-CNF formula of size $2^\ell$ as follows: $\bigwedge_{\alpha\in\{0,1\}^\ell}(s_i\ne\alpha\vee s_j\ne\alpha)$; notice that the predicate $(s_i\ne\alpha\vee s_j\ne\alpha)$ can be represented by a clause with $2\ell$ literals. If $m>2^\ell$, then the formula $\mathrm{BPHP}_{2^\ell}^m$ is unsatisfiable.
Let $\phi$ be a CNF formula with $n$ variables, and $g:\{0,1\}^k\to\{0,1\}$ be a Boolean function. Then $\phi\circ g$ denotes a CNF formula on $kn$ variables that represents $\phi(g(\vec x_1),g(\vec x_2),\ldots,g(\vec x_n))$, where $\vec x_i$ denotes a vector of $k$ new variables. $\phi\circ g$ is constructed by applying the substitution to every clause $C$ of $\phi$ and converting the resulting function $C\circ g$ to CNF in some fixed way.

% Faithful formal-publication transcription: complete Section3.
\section*{3 Communication protocols from tree-like Res(PCd) proofs}
Let $\phi$ be an unsatisfiable CNF formula with $n$ variables. Let us define the communication problem $\oplus_k\operatorname{Search}(\phi)$ with $k$ parties as follows. Assume that the $i$th party has an assignment $\alpha_i\in\{0,1\}^n$ written on the forehead. They aim to find a clause of $\phi$ falsified by the assignment $\sum_{i=1}^k\alpha_i$ (all sums of boolean vectors are computed modulo 2).
\begin{authorclaim}{Lemma1}
Let $\phi$ be an unsatisfiable CNF formula. If there exists a tree-like $\mathrm{Res}(\mathrm{PC}_d)$ proof of $\phi$ of length $m$, then $R^{1/3}_{\mathrm{pub}}(\oplus_{d+1}\operatorname{Search}(\phi))=O(d\cdot\log m)$.
\end{authorclaim}
A slightly weaker version of the following lemma was implicitly proved in [15]:
\begin{authorclaim}{Lemma12 (see proof of Theorem3.11 from [15])}
Let $T$ be a binary tree with $m$ vertices such that the $i$th vertex is labeled with $a_i\in\{0,1\}$ with the hereditary property: for each inner vertex $i$ with direct descendants $c_1$ and $c_2$, if $a_i=1$, then $a_{c_1}=1$ or $a_{c_2}=1$. We also assume that if $r$ is the root of $T$, then $a_r=1$. Assume that we have a one-sided bounded error oracle access to $a_i$ i.e. if we request a value of $a_i$ and $a_i=0$ we get 1 with probability at most $1/2$ and0 with probability at least $1/2$; if $a_i=1$ we get 1 with probability 1. Then there exists an algorithm $A$ that with probability at least $2/3$ returns a leaf $\ell$ of $T$ with $a_\ell=1$ and makes $O(\log m)$ oracle queries to $a_1,\ldots,a_m$.
\end{authorclaim}
Proof. See AppendixA.
\begin{proof}[Proof of Lemma1]
Let $F_1,\ldots,F_m$ be a tree-like $\mathrm{Res}(\mathrm{PC}_d)$-refutation of $\phi$ with the underlying tree $T$, where vertices of $T$ are identified with $[m]$. Then the leaves of $T$ correspond to the clauses of $\phi$ and $m$ is the root of $T$.
Let $\alpha_1,\ldots,\alpha_{d+1}$ be the assignments written on the foreheads of $d+1$ parties. Let $\alpha=\sum_{i=1}^{d+1}\alpha_i$. Let $a_i=1$ iff $\alpha$ falsifies $F_i$ for $i\in[m]$. Then $a_m=1$ since $F_m$ is identically false. For any inner node $v$ of $T$, if $a_v=1$ then for the direct descendants of $v$, $c_1$ and $c_2$ either $a_{c_1}=1$ or $a_{c_2}=1$. In the next paragraphs we show that for any $i\in[m]$ there exists a NOF $(d+1)$-party protocol that computes $a_i$ given that for each $j\in[d+1]$ the $j$th party has $\alpha_j$ written on their forehead such that
\begin{itemize}
\item the protocol transmits $O(d)$ bits;
\item the protocol has one-sided bounded error: if $a_i=1$ then the protocol returns 1 with probability 1 and if $a_i=0$ the protocol returns 0 with probability at least $1/2$.
\end{itemize}
Then we use this protocol to compute $a_i$ as an oracle in the algorithm given by Lemma12 and thus show that there is a NOF $(d+1)$-party protocol computing $\oplus_{d+1}\operatorname{Search}(\phi)$ with communication cost $O(d\log m)$.
Now we show that for every $\ell\in[m]$, $F_\ell(\alpha)$ can be computed by a $(d+1)$-party NOF protocol with one-sided error using $O(d)$ bits of communication. Let $F_\ell=\bigvee_{j=1}^t(f_j=1)$, where $f_1,\ldots,f_t$ are polynomials over $\mathbb F_2$ of degree at most $d$. Let $z_1,\ldots,z_n$ be the variables of $\phi$. Let us introduce new variables $y_{1,1},\ldots,y_{1,n},\ldots,y_{d+1,1},\ldots,y_{d+1,n}$ and assume that for each $i\in[d+1]$ the $i$th party has the value of variables $y_{i,1},y_{i,2},\ldots,y_{i,n}$ written on the forehead or in other words $\alpha_i$ assigns values of $y_{i,1},y_{i,2},\ldots,y_{i,n}$. Let $\bar f_j$ denote $f_j$ after substitution $z_\ell:=y_{1,\ell}+y_{2,\ell}+\cdots+y_{d+1,\ell}$ for $\ell\in[n]$; $j\in[t]$. Since for all $j\in[t]$, $\deg\bar f_j=\deg f_j\le d$, we can represent $\bar f_j=\bar f_j^{(1)}+\cdots+\bar f_j^{(d+1)}$ such that $\bar f_j^{(s)}$ does not contain variables $y_{s,1},\ldots,y_{s,n}$ for each $s\in[d+1]$. Then the $i$th party can compute $\bar f_1^{(i)}(\alpha_1,\ldots,\alpha_{d+1}),\ldots,\bar f_t^{(i)}(\alpha_1,\ldots,\alpha_{d+1})$. Notice that $F_\ell=\neg(\bigwedge_{j=1}^t(f_j=0))$.
The final step of the protocol exploits the idea used to construct a short randomized communication protocol for equality. Take a random uniformly distributed vector $(e_1,\ldots,e_t)\in\mathbb F_2^t$. Then all parties compute
\[\sum_{j=1}^te_jf_j(\alpha)=\sum_{i=1}^{d+1}\underbrace{\sum_{j=1}^te_j\bar f_j^{(i)}}_{i\text{th party}}\]
with $O(d)$ bits of communication and the protocol halts. To bound the error probability we use the following well-known statement:
\begin{authorclaim}{Proposition13 (Random subsum principle)}
For any $x\in\mathbb F_2^k\setminus\{0^k\}$,
\[\Pr_{y\leftarrow U(\mathbb F_2^k)}\left[\sum_{i=1}^ky_ix_i=1\right]=\frac12.\]
\end{authorclaim}
If $F_\ell(\alpha)=1$ then $\Pr[\sum_{j=1}^te_jf_j(\alpha)\ne0]=1/2$ by the random subsum principle. If $F_\ell(\alpha)=0$, then $\Pr[\sum_{j=1}^te_jf_j(\alpha)=0]=1$.
\end{proof}
\begin{authorclaim}{Remark14}
Similarly to the proof of Lemma1 one can prove that if an unsatisfiable CNF formula $\phi$ has a tree-like $\mathrm T^{\mathrm{cc}}_{\mathrm{os}}(k,c)$ refutation of length $\ell$, then for any $k$-partition of the variables, there is a randomized bounded-error $k$-party NOF protocol for Search$(\phi)$ with communication cost $O(c\log\ell)$. Thus, the bound from Lemma9 can be slightly improved in the case of one-sided error.
\end{authorclaim}

In this section we prove the following theorem:
\begin{authorclaim}{Theorem 2}
The size of any tree-like semantic $\mathrm{Res}(\oplus)$ refutation of the formula $\mathrm{PM}_{K_{n+2,n}}$ is $2^{\Omega(n)}$.

\end{authorclaim}
By Lemma 1, to prove Theorem 2 it is sufficient to show that $R^{1/3}_{\mathrm{pub}}(\oplus_2\operatorname{Search}(\mathrm{PM}_{K_{n+2,n}}))=\Omega(n)$.
Consider the communication problem $\oplus\mathrm{PM}_n^m$ that is defined as follows: Alice and Bob have matrices $X$ and $Y$ over $\mathbb F_2$ respectively, each of the matrices has size $m\times n$, where $m\ne n$. Their goal is to find an all-zero row or column or two 1-cells in the same row or column in the matrix $X+Y$.
\begin{authorclaim}{Proposition 15}
$R^{1/3}_{\mathrm{pub}}(\oplus_2\operatorname{Search}(\mathrm{PM}_{K_{n+2,n}}))\ge R^{1/3}_{\mathrm{pub}}(\oplus\mathrm{PM}_n^{n+2})$.
\end{authorclaim}
\begin{proof}
A Boolean matrix of size $(n+2)\times n$ naturally corresponds to a subset of edges of $K_{n+2,n}$. A falsified clause encoding that a vertex must be covered by a matching corresponds to an all-zero row or column of the matrix; a falsified clause, encoding that a vertex can not be covered by a matching twice, corresponds to two ones in the same row or column.
\end{proof}
Theorem 2 follows from Proposition 15 and the following theorem.
\begin{authorclaim}{Theorem 16}
$R^{1/3}_{\mathrm{pub}}(\oplus\mathrm{PM}_n^{n+2})=\Omega(n)$.
\end{authorclaim}
\begin{proof}
We assume that $n=4m+1$, where $m$ is a non-negative integer. If the theorem is true for all $n$ with the residue 1 modulo 4, then it also holds for all other $n$. Indeed, the protocol for $\oplus\mathrm{PM}_{n+1}^{n+3}$ can be used for $\oplus\mathrm{PM}_n^{n+2}$ by adding to Alice's matrix an extra column and a row with exactly one 1-cell on their intersection and to Bob's matrix an extra column and a row with all zeros.
Let $P_0$ be a protocol for $\oplus\mathrm{PM}_n^{n+2}$ transmitting at most $k$ bits. We are going to apply $P_0(X,Y)$ only to the instances where the matrix $X+Y$ does not contain all-zero rows or columns. Thus, we assume that with probability at least $2/3$ $P_0$ returns a tuple $(r_1,c_1,r_2,c_2)\in([n+2]\times[n])^2$ such that $(X+Y)_{r_1,c_1}=(X+Y)_{r_2,c_2}=1$ and either
\[
\begin{cases}r_1=r_2\\c_1\ne c_2\end{cases}
\quad\text{or}\quad
\begin{cases}r_1\ne r_2\\c_1=c_2.\end{cases}
\]
With $O(1)$ bits of communication Alice and Bob can verify whether the answer of $P_0$ is correct and return $\bot$ (failure) if it is not. Also, we can reduce the failure probability by the repetition of the protocol. Let $P$ be a protocol for $\oplus\mathrm{PM}_n^{n+2}$ under the promise that $X+Y$ does not contain all-zero rows and columns that uses $O(k)$ bits of communication and returns a correct answer with probability at least $99/100$ and $\bot$ otherwise.
We are going to construct a protocol for $\mathrm{DISJ}_m$ transmitting $O(k)$ bits, where $m=(n-1)/4$. Since by Theorem 8 any protocol for $\mathrm{DISJ}_m$ transmits $\Omega(m)$ bits, we conclude that $k=\Omega(m)$. Let Alice's input for $\mathrm{DISJ}_m$ be $a_1,\ldots,a_m$ and Bob's input be $b_1,\ldots,b_m$.
\begin{authorclaim}{Lemma 17}
There exist matrices $A(0),A(1),B(0),B(1)\in\mathbb F_2^{4\times4}$ such that $A(x)+B(y)$ is a permutation matrix iff $x\wedge y$ is 0 and
\begin{equation}
A(1)+B(1)=\begin{pmatrix}0&1&1&1\\1&0&0&0\\1&0&0&0\\1&0&0&0\end{pmatrix}.
\tag{1}
\end{equation}
\end{authorclaim}
\begin{proof}
We simply present matrices that satisfy the conditions:
\[
A(0)=\begin{pmatrix}0&1&1&0\\1&1&0&0\\1&0&1&0\\0&0&0&0\end{pmatrix};
\quad A(1)=\begin{pmatrix}0&0&0&0\\0&0&0&0\\0&0&0&0\\0&0&0&0\end{pmatrix};
\]
\[
B(0)=\begin{pmatrix}0&0&1&0\\0&1&0&0\\1&0&0&0\\0&0&0&1\end{pmatrix};
\quad B(1)=\begin{pmatrix}0&1&1&1\\1&0&0&0\\1&0&0&0\\1&0&0&0\end{pmatrix}.
\]
\end{proof}
Notice that Lemma 17 immediately allows to reduce $\mathrm{DISJ}_m$ to the problem of checking whether the sum of Alices and Bobs matrices is a permutation matrix. In order to achieve that, Alice builds a matrix $A=\operatorname{diag}(A(a_1),\ldots,A(a_m))$, Bob builds a matrix $B=\operatorname{diag}(B(b_1),\ldots,B(b_m))$. It is easy to see that $A+B$ is a permutation matrix iff $\mathrm{DISJ}_m(a,b)=1$.
Let us describe the reduction of $\mathrm{DISJ}_m$ to $\oplus\mathrm{PM}_n^{n+2}$. Alice and Bob first construct matrices $X_0$ and $Y_0$ of the following form:
\[
X_0=\begin{pmatrix}A&0_{(n-1)\times1}\\0_{1\times(n-1)}&1\\0_{1\times(n-1)}&1\\0_{1\times(n-1)}&1\end{pmatrix};
\quad Y_0=\begin{pmatrix}B&0_{(n-1)\times1}\\0_{1\times(n-1)}&0\\0_{1\times(n-1)}&0\\0_{1\times(n-1)}&0\end{pmatrix},
\]
then
\[
X_0+Y_0=\begin{pmatrix}A+B&0_{(n-1)\times1}\\0_{1\times(n-1)}&1\\0_{1\times(n-1)}&1\\0_{1\times(n-1)}&1\end{pmatrix},
\]
where $A+B$ is a permutation matrix iff $\mathrm{DISJ}_m(a,b)=1$. Then if $P(X_0,Y_0)$ returns two cells that do not belong to the column $n$ we may conclude that $\mathrm{DISJ}_m(a,b)=0$. If $P(X_0,Y_0)$ returns two cells from the $n$th column, then the value of $\mathrm{DISJ}_m(a,b)$ can not be uniquely determined. Notice that for $X_0$ and $Y_0$ constructed as above the protocol always returning $(n+1,n,n+2,n)$ solves $\oplus\mathrm{PM}_n^{n+2}$.
If $\mathrm{DISJ}_m(a,b)=0$, then the matrix $X_0+Y_0$ contains at least two columns with three ones and these columns are indistinguishable from each over. To make use of that, we randomly shuffle rows and columns.
We are going to construct a protocol $T$ for $\mathrm{DISJ}_m$ as follows: Alice and Bob choose permutations $\pi\in S_n$, $\tau\in S_{n+2}$ and a matrix $\Delta\in\mathbb F_2^{(n+2)\times n}$ uniformly at random. We define matrices $X_0^{\tau,\pi}$ and $Y_0^{\tau,\pi}$ from $\mathbb F_2^{(n+2)\times n}$ such that for each $i\in[n+2]$ and $j\in[n]$, $(X_0^{\tau,\pi})_{i,j}=(X_0)_{\tau(i),\pi(j)}$ and $(Y_0^{\tau,\pi})_{i,j}=(Y_0)_{\tau(i),\pi(j)}$. Alice and Bob run the protocol $P$ for inputs $X=X_0^{\tau,\pi}+\Delta$, $Y=Y_0^{\tau,\pi}+\Delta$. Notice that $X+Y=X_0^{\tau,\pi}+Y_0^{\tau,\pi}$, thus $X+Y$ can be obtained from $X_0+Y_0$ by shuffling rows and columns. If $P(X,Y)$ returns two cells from the column $\pi(n)$, Alice and Bob return 1, if $P(X,Y)$ returns two cells from other column or row, Alice and Bob return 0. If $P(X,Y)$ returns $\bot$, then Alice and Bob return $\bot$.
First notice that if $\mathrm{DISJ}_m(a,b)=1$, then $T$ returns a correct answer or $\bot$ with probability 1 (and the probability of $\bot$ is at most $1/100$), since in that case $X+Y$ has exactly one column with three 1-cells, each of the other columns and rows contains exactly one 1-cell. Let us fix $a,b\in\{0,1\}^m$ such that $\mathrm{DISJ}_m(a,b)=0$. We denote $p:=\Pr[T(a,b)=0]$, we will show that $p\ge99/200$. We can then increase this probability to $2/3$ by repeating the protocol twice (if $T(a,b)$ returns 0 at least once, we return 0, if $T(a,b)$ always return $\bot$, we return $\bot$, otherwise we return 1).
Let us describe random bits used by the constructed protocol $T$. First, we use random bits $r$ to run the protocol $P$. Second, we use random bits to generate $\pi,\tau$, and $\Delta$. Since $\mathrm{DISJ}_m(a,b)=0$, we can fix $i\in[m]$ such that $a_i=b_i=1$. In that case the submatrix of $X_0+Y_0$ formed by rows and columns with the indices $4(i-1)+1,4(i-1)+2,4(i-1)+3,4(i-1)+4$ coincides with the matrix (1). Let us denote by $\operatorname{col}(j)$ for $j\in[n]$ the set of all tuples $(x,j,y,j)\in([n+2]\times[n])^2$.
\begin{align*}
p&=\Pr_{\pi,\tau,\Delta,r}[P_r(X,Y)\notin\operatorname{col}(\pi(n))]
-\underbrace{\Pr_{\pi,\tau,\Delta,r}[P_r(X,Y)=\bot]}_{=:p_\bot}\\
&=1-\Pr_{\pi,\tau,\Delta,r}[P_r(X,Y)\in\operatorname{col}(\pi(n))]-p_\bot\\
&=1-\sum_{\pi_0,\tau_0}\Pr_{r,\Delta}[P_r(X_0^{\tau_0,\pi_0}+\Delta,Y_0^{\tau_0,\pi_0}+\Delta)\in\operatorname{col}(\pi_0(n))]
\Pr_{\pi,\tau}[\pi=\pi_0,\tau=\tau_0]-p_\bot .
\end{align*}
Observe that for fixed $\pi_0$ and $\tau_0$ the random variable $(X_0^{\tau_0,\pi_0}+\Delta,Y_0^{\tau_0,\pi_0}+\Delta)$ is uniformly distributed over the pairs of matrices with the sum $X_0^{\tau_0,\pi_0}+Y_0^{\tau_0,\pi_0}$. Let $\alpha\in S_n$ be the transposition swapping $n$ and $4(i-1)+1$. Let $\beta\in S_{n+2}$ be the permutation swapping $n$ and $4(i-1)+2$, $n+1$ and $4(i-1)+3$, $n+2$ and $4(i-1)+3$ (i.e. $\beta$ is a product of three transpositions). By the construction of $\alpha$ and $\beta$, $(X_0+Y_0)=(X_0^{\beta,\alpha}+Y_0^{\beta,\alpha})$, thus $(X_0^{\tau,\pi}+Y_0^{\tau,\pi})=(X_0^{\tau\circ\beta,\pi\circ\alpha}+Y_0^{\tau\circ\beta,\pi\circ\alpha})$ for every $\pi,\tau$. Thus the random variable $(X_0^{\tau_0\circ\beta,\pi_0\circ\alpha}+\Delta,Y_0^{\tau_0\circ\beta,\pi_0\circ\alpha}+\Delta)$ has the same distribution with $(X_0^{\tau_0,\pi_0}+\Delta,Y_0^{\tau_0,\pi_0}+\Delta)$, thus we can continue the sequence as follows:
\begin{align*}
p&=1-\sum_{\pi_0,\tau_0}\Pr_{r,\Delta}[P_r(X_0^{\tau_0\circ\beta,\pi_0\circ\alpha}+\Delta,Y_0^{\tau_0\circ\beta,\pi_0\circ\alpha}+\Delta)\in\operatorname{col}(\pi_0(n))]\Pr_{\pi,\tau}[\pi=\pi_0,\tau=\tau_0]-p_\bot\\
&=1-\sum_{\pi_0,\tau_0}\Pr_{r,\Delta}[P_r(X_0^{\tau_0,\pi_0}+\Delta,Y_0^{\tau_0,\pi_0}+\Delta)\in\operatorname{col}(\pi_0\circ\alpha^{-1}(n))]\Pr_{\pi,\tau}[\pi=\pi_0,\tau=\tau_0]-p_\bot\\
&=1-\Pr_{\pi,\tau,\Delta,r}[P_r(X,Y)\in\operatorname{col}((\pi\circ\alpha^{-1})(n))]-p_\bot\\
&\ge1-\Pr_{\pi,\tau,\Delta,r}[P_r(X,Y)\notin\operatorname{col}(\pi(n))]-p_\bot=1-p-p_\bot .
\end{align*}
Thus, $p\ge1-p-p_\bot$ and $p\ge(1-p_\bot)/2=99/200$.
\end{proof}

\section*{A Proof of Lemma 12}
\begin{authorclaim}{Lemma 12}
Let $T$ be a binary tree with $m$ vertices such that the $i$th vertex is labeled with $a_i\in\{0,1\}$ with the hereditary property: for each inner vertex $i$ with direct descendants $c_1$ and $c_2$, if $a_i=1$, then $a_{c_1}=1$ or $a_{c_2}=1$. We also assume that if $r$ is the root of $T$, then $a_r=1$. Assume that we have a one-sided bounded error oracle access to $a_i$ i.e. if we request a value of $a_i$ and $a_i=0$ we get 1 with probability at most $1/2$ and 0 with probability at least $1/2$; if $a_i=1$ we get 1 with probability 1. Then there exists an algorithm $A$ that with probability at least $2/3$ returns a leaf $\ell$ of $T$ with $a_\ell=1$ and makes $O(\log m)$ oracle queries to $a_1,\ldots,a_m$.
\end{authorclaim}
\begin{proof}[Proof of Lemma 12]
For a tree $F$ we denote by $|F|$ the number of nodes in $F$ and for a node $v$ of $F$ we denote by $\operatorname{Subtree}(F,v)$ the subtree of $F$ with root $v$. Let Oracle$(i)$ be the oracle function returning the correct value of $a_i$ with probability at least $9/10$. We can implement such a function using the majority vote of a constant number of initial oracle queries. Let $C$ be a constant; an appropriate value of $C$ we choose later. Consider Algorithm 2 on the following page.
We claim that at any iteration $T_i$ has the hereditary property. This is the case in the beginning and if $i$ decreases at some iteration, then the next $T_i$ was considered at an earlier iteration. Otherwise, the next $T_i$ is either a subtree of the current $T_i$ (in that case the hereditary property is clearly maintained), or is obtained by removal a subtree with 0-labeled root (here we use that the oracle has a one-sided error) from the previous $T_i$ (the hereditary property is also maintained in that case).
We first consider a variant of the algorithm that works infinitely long (i.e., $C=+\infty$) and compute the expected number of the first iteration such that $T_i$ consists of a single 1-labeled leaf of $T$. Notice that after the first such iteration the value of $T_i$ stays the same for all further iterations. We show that that the expected value is at most $C\log m$ for some constant $C$. Then by running the algorithm for $3C\lceil\log m\rceil$ iterations we obtain the required error probability by Markov's inequality.
\begin{authorclaim}{Algorithm 2: Search for 1-leaf}
\begin{algorithmic}
\State $T_0:=T$ \Comment{Initialize the tree}
\State $i:=0$
\For{$j:=1$ to $3C\lceil\log_{3/2}m\rceil$}
 \State $r:=$ root of $T_i$
 \If{Oracle$(r)=0$}
  \State $i:=\max\{0,i-1\}$ \Comment{Backtrack since the current tree may not contain a 1-leaf}
 \ElsIf{$|T_i|\ne1$}
  \State $v:=$ a centroid node of $T_i$
  \Statex \hspace{\algorithmicindent} $\triangleright$ i.e. such that $|\operatorname{Subtree}(T_i,v)|\in[\frac13|T_i|,\frac23|T_i|]$
  \If{Oracle$(v)=1$}
   \State $T_{i+1}:=\operatorname{Subtree}(T_i,v)$
  \Else
   \State $T_{i+1}:=T_i-\operatorname{Subtree}(T_i,v)$
   \Statex \hspace{\algorithmicindent} $\triangleright$ $T_{i+1}$ is obtained from $T_i$ by the deletion of $\operatorname{Subtree}(T_i,v)$
  \EndIf
  \State $i:=i+1$
 \EndIf
\EndFor
\State \Return the only node of $T_i$, if $|T_i|=1$
\end{algorithmic}
\end{authorclaim}
Let $T(j)$ denote the value of $T_i$ before the start of $j$th iteration, $i(j)$ denote $i$ at the start of $j$th iteration and $r(j)$ denote the root of $T(j)$. Notice that if $a_{r(j)}=1$, then for every $j'>j$, $T(j')$ is a subtree of $T(j)$, since the algorithm never backtracks if the true value of the roots label is 1. Hence, if $a_{r(j)}=a_{r(j')}=1$ for some $j<j'$, then $i(j)\le i(j')$.
Let us consider a sequence $j_1,j_2,j_3,\ldots$, where $j_1=0$, $j_s=\min\{j\mid a_{r(j)}=1\wedge j>j_{s-1}\wedge i(j)>i(j_{s-1})\}$, if such minimum exists.
Let us consider the iterations from $j_s$ till $j_{s+1}-1$. We consider the random variables $Y_{j_s},Y_{j_s+1},\ldots,Y_{j_{s+1}-1}$ corresponding to these iterations with the following properties:
\begin{itemize}
\item If $T(j)$ coincides with $T(j_s)$, then its root is labeled with 1. Then $Y_j=-1$ if the second oracle query returns the correct answer and $Y_j=1$ if the answer it incorrect. Notice that $\Pr[Y_j=-1]\ge9/10$.
\item If the root of $T(j)$ is labeled with zero, then $Y_j=-1$, if the first oracle query returns the correct answer (i.e. the algorithm backtracks). Otherwise, if $T(j)$ consists of a single node $Y_j=0$. Otherwise, if the root of $T(j+1)$ is labeled with 0, then $Y_j=1$. If it is labeled with 1, then $Y_j=-\infty$. Notice that $\Pr[Y_j\le-1]\ge9/10$.
\end{itemize}
Notice that, $j_{s+1}=j_s+\min\{k\mid\sum_{j=j_s}^{j_s+k-1}Y_j\le-1\}$. In order to estimate the expected value of $j_{s+1}-j_s$ we consider an auxiliary random variables $X_{j_s},X_{j_s+1},\ldots,X_{j_{s+1}-1}$, defined as
\[
X_j=\begin{cases}1,&\text{if }Y_j\ge0\\-1,&\text{if }Y_j<0.\end{cases}
\]
Notice then $\sum_{j=j_s}^{j_s+k-1}Y_j\le\sum_{j=j_s}^{j_s+k-1}X_j$. We can apply the following fact about random walks in a straight line to the random variables $X_j$:
\begin{authorclaim}{Theorem 45 (Section XII.2 of [18])}
Let $X_1,X_2,\ldots$ be a sequence of independent random variables that take value in $\{-1,1\}$. Assume that for all $i$, $\Pr[X_i=1]\le1/10$ and $\Pr[X_i=-1]\ge9/10$. Let $M$ be a random variable that equals the minimal natural number $k$ such that $\sum_{i=1}^kX_i=-1$. Then the expected value of $M$ is at most $C$, where $C\in\mathbb R$ is an absolute constant.
\end{authorclaim}
Fact 45 implies that $\mathbb E[j_{s+1}-j_s]\le C$. Then $\mathbb E[j_s]=\mathbb E[j_s-j_{s-1}+(j_{s-1}-j_{s-2})+\cdots+(j_2-j_1)+(j_1-j_0)]\le sC$. Thus, by Markov's inequality $\Pr[j_s\le3sC]\ge2/3$. Since $|T_{j_s}|\le(2/3)^s|T_{j_0}|$, the algorithm that runs for $3C\lceil\log_{3/2}m\rceil$ iterations terminates in a 1-labeled leaf with probability at least $2/3$.
\end{proof}

\begin{thebibliography}{33}
\bibitem{r1} Paul Beame, Toniann Pitassi, and Nathan Segerlind. Lower bounds for Lovász-Schrijver systems and beyond follow from multiparty communication complexity. SIAM J. Comput., 37(3):845–869, 2007.
\bibitem{r2} Paul Beame and Søren Riis. More on the relative strength of counting principles. In Paul Beam and Samuel R. Buss, editors, Proof Complexity and Feasible Arithmetics, Proceedings of a DIMACS Workshop, New Brunswick, New Jersey, USA, April 21-24, 1996, volume 39 of DIMACS Series in Discrete Mathematics and Theoretical Computer Science, pages 13–35. DIMACS/AMS, 1996. doi:10.1090/dimacs/039/02.
\bibitem{r3} Stephen A. Cook and Robert A. Reckhow. The relative efficiency of propositional proof systems. The Journal of Symbolic Logic, 44(1):36–50, 1979.
\bibitem{r4} William Cook, Collette R. Coullard, and Gy. Turán. On the complexity of cutting-plane proofs. Discrete Applied Mathematics, 18(1):25–38, 1987.
\bibitem{r5} Stefan S. Dantchev, Nicola Galesi, and Barnaby Martin. Resolution and the binary encoding of combinatorial principles. In Amir Shpilka, editor, 34th Computational Complexity Conference, CCC 2019, July 18-20, 2019, New Brunswick, NJ, USA, volume 137 of LIPIcs, pages 6:1–6:25. Schloss Dagstuhl - Leibniz-Zentrum für Informatik, 2019. doi:10.4230/LIPIcs.CCC.2019.6.
\bibitem{r6} Noah Fleming, Denis Pankratov, Toniann Pitassi, and Robert Robere. Random $\Theta(\log n)$CNFs are hard for cutting planes. In Chris Umans, editor, 58th IEEE Annual Symposium on Foundations of Computer Science, FOCS 2017, Berkeley, CA, USA, October 15-17, 2017, pages 109–120. IEEE Computer Society, 2017. doi:10.1109/FOCS.2017.19.
\bibitem{r7} Ankit Garg, Mika Göös, Pritish Kamath, and Dmitry Sokolov. Monotone circuit lower bounds from resolution. In Ilias Diakonikolas, David Kempe, and Monika Henzinger, editors, Proceedings of the 50th Annual ACM SIGACT Symposium on Theory of Computing, STOC 2018, Los Angeles, CA, USA, June 25-29, 2018, pages 902–911. ACM, 2018. doi:10.1145/ 3188745.3188838.
\bibitem{r8} Michal Garlík and Leszek Aleksander Kolodziejczyk. Some subsystems of constant-depth frege with parity. ACM Trans. Comput. Log., 19(4):29:1–29:34, 2018. doi:10.1145/3243126.
\bibitem{r9} Mika Göös and Toniann Pitassi. Communication lower bounds via critical block sensitivity. In Proceedings of the Forty-Sixth Annual ACM Symposium on Theory of Computing, STOC ’14, page 847–856, New York, NY, USA, 2014. Association for Computing Machinery. doi: 10.1145/2591796.2591838.
\bibitem{r10} Pavel Hrubes and Pavel Pudlák. Random formulas, monotone circuits, and interpolation. In Chris Umans, editor, 58th IEEE Annual Symposium on Foundations of Computer Science, FOCS 2017, Berkeley, CA, USA, October 15-17, 2017, pages 121–131. IEEE Computer Society, 2017. doi:10.1109/FOCS.2017.20.
\bibitem{r11} Pavel Hrubes and Pavel Pudlák. A note on monotone real circuits. Inf. Process. Lett., 131:15–19, 2018. doi:10.1016/j.ipl.2017.11.002.
\bibitem{r12} Trinh Huynh and Jakob Nordström. On the virtue of succinct proofs: amplifying communication complexity hardness to time-space trade-offs in proof complexity. In Howard J. Karloff and Toniann Pitassi, editors, Proceedings of the 44th Symposium on Theory of Computing Conference, STOC 2012, New York, NY, USA, May 19 - 22, 2012, pages 233–248. ACM, 2012. doi:10.1145/2213977.2214000.
\bibitem{r13} Russell Impagliazzo, Toniann Pitassi, and Alasdair Urquhart. Upper and lower bounds for tree-like cutting planes proofs. In Proceedings of the Ninth Annual Symposium on Logic in Computer Science (LICS ’94), Paris, France, July 4-7, 1994, pages 220–228. IEEE Computer Society, 1994. doi:10.1109/LICS.1994.316069.
\bibitem{r14} Dmitry Itsykson and Dmitry Sokolov. Lower bounds for splittings by linear combinations. In Erzsébet Csuhaj-Varjú, Martin Dietzfelbinger, and Zoltán Ésik, editors, Mathematical Foundations of Computer Science 2014 - 39th International Symposium, MFCS 2014, Budapest, Hungary, August 25-29, 2014. Proceedings, Part II, volume 8635 of Lecture Notes in Computer Science, pages 372–383. Springer, 2014. doi:10.1007/978-3-662-44465-8\_32.
\bibitem{r15} Dmitry Itsykson and Dmitry Sokolov. Resolution over linear equations modulo two. Annals of Pure and Applied Logic, 171(1), January 2020. doi:10.1016/j.apal.2019.102722.
\bibitem{r16} Bala Kalyanasundaram and Georg Schnitger. The probabilistic communication complexity of set intersection. SIAM J. Discret. Math., 5(4):545–557, 1992. doi:10.1137/0405044.
\bibitem{r17} Mauricio Karchmer and Avi Wigderson. Monotone circuits for connectivity require superlogarithmic depth. SIAM J. Discret. Math., 3(2):255–265, 1990. doi:10.1137/0403021.
\bibitem{r18} H. Kesten. An introduction to probability theory and its applications, volume i, (william feller). SIAM Review, 11(1):96–96, 1969. doi:10.1137/1011021.
\bibitem{r19} Erfan Khaniki. On proof complexity of resolution over polynomial calculus. Electronic Colloquium on Computational Complexity (ECCC), 27:34, 2020. URL: https://eccc.weizmann. ac.il/report/2020/034.
\bibitem{r20} Jan Krajíček. Interpolation theorems, lower bounds for proof systems, and independence results for bounded arithmetic. J. Symb. Log., 62(2):457–486, 1997. doi:10.2307/2275541.
\bibitem{r21} Jan Krajíček. An exponential lower bound for a constraint propagation proof system based on ordered binary decision diagrams. J. Symb. Log., 73(1):227–237, 2008. doi:10.2178/jsl/ 1208358751.
\bibitem{r22} Jan Krajíček. Randomized feasible interpolation and monotone circuits with a local oracle. J. Mathematical Logic, 18(2):1850012:1–1850012:27, 2018. doi:10.1142/S0219061318500125.
\bibitem{r23} Jan Krajíček. Proof complexity, volume 170. Cambridge University Press, 2019.
\bibitem{r24} Eyal Kushilevitz and Noam Nisan. Communication complexity. Cambridge University Press, 1997.
\bibitem{r25} Edward I Nechiporuk. A boolean function. Engl. transl. in Sov. Phys. Dokl., 10:591–593, 1966.
\bibitem{r26} Vsevolod Oparin. Tight upper bound on splitting by linear combinations for pigeonhole principle. In Nadia Creignou and Daniel Le Berre, editors, Theory and Applications of Satisfiability Testing - SAT 2016 - 19th International Conference, Bordeaux, France, July 5-8, 2016, Proceedings, volume 9710 of Lecture Notes in Computer Science, pages 77–84. Springer, 2016. doi:10.1007/978-3-319-40970-2\_6.
\bibitem{r27} Fedor Part and Iddo Tzameret. Resolution with counting: Dag-like lower bounds and different moduli. In Thomas Vidick, editor, 11th Innovations in Theoretical Computer Science Conference, ITCS 2020, January 12-14, 2020, Seattle, Washington, USA, volume 151 of LIPIcs, pages 19:1–19:37. Schloss Dagstuhl - Leibniz-Zentrum für Informatik, 2020. doi:10.4230/LIPIcs.ITCS.2020.19.
\bibitem{r28} Pavel Pudlák. Lower bounds for resolution and cutting plane proofs and monotone computations. J. Symb. Log., 62(3):981–998, 1997. doi:10.2307/2275583.
\bibitem{r29} Ran Raz and Pierre McKenzie. Separation of the monotone NC hierarchy. Combinatorica, 19(3):403–435, 1999. doi:10.1007/s004930050062.
\bibitem{r30} Ran Raz and Avi Wigderson. Monotone circuits for matching require linear depth. J. ACM, 39(3):736–744, 1992. doi:10.1145/146637.146684.
\bibitem{r31} Alexander A. Sherstov. The multiparty communication complexity of set disjointness. In Proceedings of the forty-fourth annual ACM symposium on Theory of computing, pages 525–548, 2012.
\bibitem{r32} Alexander A. Sherstov. Communication lower bounds using directional derivatives. J. ACM, 61(6), December 2014. doi:10.1145/2629334.
\bibitem{r33} Dmitry Sokolov. Dag-like communication and its applications. In Computer Science Theory and Applications - 12th International Computer Science Symposium in Russia, CSR 2017, Kazan, Russia, June 8-12, 2017, Proceedings, pages 294–307, 2017. doi:10.1007/ 978-3-319-58747-9\_26.
\end{thebibliography}
\end{document}
